% GENERATED FILE. Edit canonical lessons, then run npm run generate:books.
% canonical-source-hash: fnv1a64:ecc3562f89f5f9d1
% canonical-lessons: JA-C395-koshou, JA-C395-shako, JA-C395-fusagaru, JA-C395-yurumeru, JA-C395-ikisugiru, JA-R395-review-390-395

\chapter{Car Trouble and Wrong Turns}
\label{ch:ja-car-trouble-and-wrong-turns}
\hlchaptermodality{395}

\begin{chapteropening}
\textbf{By the end of this chapter:} \emph{I can say what went wrong on the road.}

Complete the last lesson of chapter 395: I can say what went wrong on the road.
\end{chapteropening}

\section[\texorpdfstring{\ja{こしょう} (koshō)}{こしょう (koshō)}]{\ja{こしょう} (koshō) — a breakdown; a fault}

\label{lesson:JA-C395-koshou}



\begin{quote}
Before the new one: say the Japanese for to descend, then the Japanese for to cross over.
\end{quote}

\subsection*{You'll want to know: \ja{こしょう}}
\textbf{\ja{こしょう}} — \emph{koshō} — \textquotedblleft{}a breakdown; a fault\textquotedblright{}.

\begin{quote}
\textbf{\ja{くるまが} \ja{こしょうしました。}} — \emph{kuruma ga koshō shimashita} — \textquotedblleft{}The car broke down.\textquotedblright{}
\end{quote}

\begin{cousinweb}
\textbf{\ja{こしょう}} (\emph{koshō}) is Sino-Japanese (kango): two Chinese characters read with their on'yomi, readings borrowed from Chinese. They are \emph{ko}, \textquotedblleft{}circumstances\textquotedblright{} (Middle Chinese \emph{kuH}) and \emph{shō}, \textquotedblleft{}hinder\textquotedblright{} (\emph{tsyangH}) (KANJIDIC2; Middle Chinese in Baxter's transcription).
\end{cousinweb}

\begin{grammarlens}[title={What you've built}]
One more word for getting around.
\end{grammarlens}

\subsection*{Your turn}
\begin{itemize}
\raggedright
  \item \emph{Say it:} \emph{koshō}
  \item \emph{Say it:} \emph{koshō}, once more
  \item \emph{Say it:} \emph{kuruma ga koshō shimashita} — The car broke down.
  \item \emph{Recall:} say the Japanese for to descend, then the Japanese for to cross over, then say \emph{koshō} again
\end{itemize}

\begin{tcolorbox}[breakable,colback=teal!4,colframe=teal!35!black,title={Before you move on}]
What does \ja{こしょう} mean? (\textquotedblleft{}A breakdown; a fault\textquotedblright{}.) Say it once more.
\end{tcolorbox}
\section[\texorpdfstring{\ja{しゃこ} (shako)}{しゃこ (shako)}]{\ja{しゃこ} (shako) — a garage; a depot}

\label{lesson:JA-C395-shako}



\begin{quote}
Before the new one: say the Japanese for to cross over, then the Japanese for a breakdown.
\end{quote}

\subsection*{You'll want to know: \ja{しゃこ}}
\textbf{\ja{しゃこ}} — \emph{shako} — \textquotedblleft{}a garage; a depot\textquotedblright{}.

\begin{quote}
\textbf{\ja{くるまを} \ja{しゃこに} \ja{いれました。}} — \emph{kuruma o shako ni iremashita} — \textquotedblleft{}I put the car in the garage.\textquotedblright{}
\end{quote}

\begin{cousinweb}
\textbf{\ja{しゃこ}} (\emph{shako}) is Sino-Japanese (kango): two Chinese characters read with their on'yomi, readings borrowed from Chinese. They are \emph{sha}, \textquotedblleft{}car\textquotedblright{} (Middle Chinese \emph{tsyhae}) and \emph{ko}, \textquotedblleft{}warehouse\textquotedblright{} (\emph{khuH}) (KANJIDIC2; Middle Chinese in Baxter's transcription).
\end{cousinweb}

\begin{grammarlens}[title={What you've built}]
One more word for getting around.
\end{grammarlens}

\subsection*{Your turn}
\begin{itemize}
\raggedright
  \item \emph{Say it:} \emph{shako}
  \item \emph{Say it:} \emph{shako}, once more
  \item \emph{Say it:} \emph{kuruma o shako ni iremashita} — I put the car in the garage.
  \item \emph{Recall:} say the Japanese for to cross over, then the Japanese for a breakdown, then say \emph{shako} again
\end{itemize}

\begin{tcolorbox}[breakable,colback=teal!4,colframe=teal!35!black,title={Before you move on}]
What does \ja{しゃこ} mean? (\textquotedblleft{}A garage; a depot\textquotedblright{}.) Say it once more.
\end{tcolorbox}
\section[\texorpdfstring{\ja{ふさがる} (fusagaru)}{ふさがる (fusagaru)}]{\ja{ふさがる} (fusagaru) — to be blocked; to be closed off}

\label{lesson:JA-C395-fusagaru}



\begin{quote}
Before the new one: say the Japanese for a breakdown, then the Japanese for a garage.
\end{quote}

\subsection*{You'll want to know: \ja{ふさがる}}
\textbf{\ja{ふさがる}} — \emph{fusagaru} — \textquotedblleft{}to be blocked; to be closed off\textquotedblright{}.

\begin{quote}
\textbf{\ja{じこで} \ja{みちが} \ja{ふさがっています。}} — \emph{jiko de michi ga fusagatte imasu} — \textquotedblleft{}The road is blocked because of an accident.\textquotedblright{}
\end{quote}

Polite \textbf{\ja{ふさがります}}, \emph{fusagarimasu}; a godan verb. Past (polite): \textbf{\ja{ゆきで} \ja{みちが} \ja{ふさがりました。}}, \emph{yuki de michi ga fusagarimashita}.

\begin{cousinweb}
\textbf{\ja{ふさがる}} (\emph{fusagaru}) is a native Japanese word (wago), inherited rather than borrowed from Chinese; its kanji was attached later as a kun'yomi, and JMdict marks it as usually written in kana (JMdict). It is the intransitive partner of \emph{fusagu}, \textquotedblleft{}to block\textquotedblright{} (JMdict).
\end{cousinweb}

\begin{grammarlens}[title={What you've built}]
One more word for getting around.
\end{grammarlens}

\subsection*{Your turn}
\begin{itemize}
\raggedright
  \item \emph{Say it:} \emph{fusagaru}
  \item \emph{Say it:} \emph{fusagaru}, once more
  \item \emph{Say it:} \emph{jiko de michi ga fusagatte imasu} — The road is blocked because of an accident.
  \item \emph{Recall:} say the Japanese for a breakdown, then the Japanese for a garage, then say \emph{fusagaru} again
\end{itemize}

\begin{tcolorbox}[breakable,colback=teal!4,colframe=teal!35!black,title={Before you move on}]
What does \ja{ふさがる} mean? (\textquotedblleft{}To be blocked; to be closed off\textquotedblright{}.) Say it once more.
\end{tcolorbox}
\section[\texorpdfstring{\ja{ゆるめる} (yurumeru)}{ゆるめる (yurumeru)}]{\ja{ゆるめる} (yurumeru) — to loosen; to ease off}

\label{lesson:JA-C395-yurumeru}



\begin{quote}
Before the new one: say the Japanese for a garage, then the Japanese for to be blocked.
\end{quote}

\subsection*{You'll want to know: \ja{ゆるめる}}
\textbf{\ja{ゆるめる}} — \emph{yurumeru} — \textquotedblleft{}to loosen; to ease off\textquotedblright{}.

\begin{quote}
\textbf{\ja{かどの} \ja{まえで} \ja{そくどを} \ja{ゆるめました。}} — \emph{kado no mae de sokudo o yurumemashita} — \textquotedblleft{}I slowed down before the corner.\textquotedblright{}
\end{quote}

Polite \textbf{\ja{ゆるめます}}, \emph{yurumemasu}; an ichidan verb. Past (polite): \textbf{\ja{うんてんしゅが} \ja{そくどを} \ja{ゆるめました。}}, \emph{untenshu ga sokudo o yurumemashita}.

\begin{cousinweb}
\textbf{\ja{ゆるめる}} (\emph{yurumeru}) is a native Japanese word (wago), inherited rather than borrowed from Chinese; its kanji was attached later as a kun'yomi (JMdict).
\end{cousinweb}

\begin{grammarlens}[title={What you've built}]
One more word for getting around.
\end{grammarlens}

\subsection*{Your turn}
\begin{itemize}
\raggedright
  \item \emph{Say it:} \emph{yurumeru}
  \item \emph{Say it:} \emph{yurumeru}, once more
  \item \emph{Say it:} \emph{kado no mae de sokudo o yurumemashita} — I slowed down before the corner.
  \item \emph{Recall:} say the Japanese for a garage, then the Japanese for to be blocked, then say \emph{yurumeru} again
\end{itemize}

\begin{tcolorbox}[breakable,colback=teal!4,colframe=teal!35!black,title={Before you move on}]
What does \ja{ゆるめる} mean? (\textquotedblleft{}To loosen; to ease off\textquotedblright{}.) Say it once more.
\end{tcolorbox}
\section[\texorpdfstring{\ja{いきすぎる} (ikisugiru)}{いきすぎる (ikisugiru)}]{\ja{いきすぎる} (ikisugiru) — to go past one's turning; to go too far}

\label{lesson:JA-C395-ikisugiru}



\begin{quote}
Before the new one: say the Japanese for to be blocked, then the Japanese for to loosen.
\end{quote}

\subsection*{You'll want to know: \ja{いきすぎる}}
\textbf{\ja{いきすぎる}} — \emph{ikisugiru} — \textquotedblleft{}to go past one's turning; to go too far\textquotedblright{}.

\begin{quote}
\textbf{\ja{しんごうを} \ja{いきすぎて、みちに} \ja{まよいました。}} — \emph{shingō o ikisugite, michi ni mayoimashita} — \textquotedblleft{}I went past the traffic light and got lost.\textquotedblright{}
\end{quote}

Polite \textbf{\ja{いきすぎます}}, \emph{ikisugimasu}; an ichidan verb. Past (polite): \textbf{\ja{ちちが} \ja{かどを} \ja{いきすぎました。}}, \emph{chichi ga kado o ikisugimashita}.

\begin{cousinweb}
\textbf{\ja{いきすぎる}} (\emph{ikisugiru}) is a native Japanese word (wago), inherited rather than borrowed from Chinese, built from \emph{iki}, \textquotedblleft{}going\textquotedblright{} (from \emph{iku}) and \emph{sugiru}, \textquotedblleft{}to go beyond\textquotedblright{}; its kanji were attached later as kun'yomi (JMdict).
\end{cousinweb}

\begin{grammarlens}[title={What you've built}]
One more word for getting around.
\end{grammarlens}

\subsection*{Your turn}
\begin{itemize}
\raggedright
  \item \emph{Say it:} \emph{ikisugiru}
  \item \emph{Say it:} \emph{ikisugiru}, once more
  \item \emph{Say it:} \emph{shingō o ikisugite, michi ni mayoimashita} — I went past the traffic light and got lost.
  \item \emph{Recall:} say the Japanese for to be blocked, then the Japanese for to loosen, then say \emph{ikisugiru} again
\end{itemize}

\begin{tcolorbox}[breakable,colback=teal!4,colframe=teal!35!black,title={Before you move on}]
What does \ja{いきすぎる} mean? (\textquotedblleft{}To go past one's turning; to go too far\textquotedblright{}.) Say it once more.
\end{tcolorbox}
\section[(dialogue)]{Looking back: from On the Expressway to Car Trouble and Wrong Turns}

\label{lesson:JA-R395-review-390-395}



\begin{quote}
Say the Japanese for an expressway and a traffic lane.
\end{quote}

\subsection*{Your turn}
Say these aloud:

\begin{itemize}
\raggedright
  \item an expressway, a traffic lane, speed, a licence, to overtake
  \item a pedestrian crossing, a pedestrian, one-way traffic, to stop walking, to block
  \item a stop on the way, a detour, a dead end, to turn back, to stop by
  \item a landmark, a signboard, a road sign, an arrow (on a sign), to look back
  \item a hill road, steep, width, to descend, to cross over
  \item a breakdown, a garage, to be blocked, to loosen, to go past one's turning
\end{itemize}

\begin{tcolorbox}[breakable,colback=teal!4,colframe=teal!35!black,title={Before you move on}]
An expressway? (\textbf{\ja{こうそくどうろ}}, \emph{kōsokudōro}.) A traffic lane? (\textbf{\ja{しゃせん}}, \emph{shasen}.) Speed? (\textbf{\ja{そくど}}, \emph{sokudo}.) A licence? (\textbf{\ja{めんきょ}}, \emph{menkyo}.) To overtake? (\textbf{\ja{おいこす}}, \emph{oikosu}.) A pedestrian crossing? (\textbf{\ja{おうだんほどう}}, \emph{ōdanhodō}.) A pedestrian? (\textbf{\ja{ほこうしゃ}}, \emph{hokōsha}.) One-way traffic? (\textbf{\ja{いっぽうつうこう}}, \emph{ippōtsūkō}.) To stop walking? (\textbf{\ja{たちどまる}}, \emph{tachidomaru}.) To block? (\textbf{\ja{さえぎる}}, \emph{saegiru}.) A stop on the way? (\textbf{\ja{よりみち}}, \emph{yorimichi}.) A detour? (\textbf{\ja{まわりみち}}, \emph{mawarimichi}.) A dead end? (\textbf{\ja{いきどまり}}, \emph{ikidomari}.) To turn back? (\textbf{\ja{ひきかえす}}, \emph{hikikaesu}.) To stop by? (\textbf{\ja{たちよる}}, \emph{tachiyoru}.) A landmark? (\textbf{\ja{めじるし}}, \emph{mejirushi}.) A signboard? (\textbf{\ja{かんばん}}, \emph{kanban}.) A road sign? (\textbf{\ja{ひょうしき}}, \emph{hyōshiki}.) An arrow (on a sign)? (\textbf{\ja{やじるし}}, \emph{yajirushi}.) To look back? (\textbf{\ja{ふりむく}}, \emph{furimuku}.) A hill road? (\textbf{\ja{さかみち}}, \emph{sakamichi}.) Steep? (\textbf{\ja{けわしい}}, \emph{kewashii}.) Width? (\textbf{\ja{はば}}, \emph{haba}.) To descend? (\textbf{\ja{くだる}}, \emph{kudaru}.) To cross over? (\textbf{\ja{こえる}}, \emph{koeru}.) A breakdown? (\textbf{\ja{こしょう}}, \emph{koshō}.) A garage? (\textbf{\ja{しゃこ}}, \emph{shako}.) To be blocked? (\textbf{\ja{ふさがる}}, \emph{fusagaru}.) To loosen? (\textbf{\ja{ゆるめる}}, \emph{yurumeru}.) To go past one's turning? (\textbf{\ja{いきすぎる}}, \emph{ikisugiru}.)
\end{tcolorbox}
